% !TEX root = ../main.tex
% ---------------------------------------------------------------------
\section{Démonstrations}
\label{sec-demos}
% ---------------------------------------------------------------------

\subsection{Propriété (1) : un décalage entier translate le spectrogramme}
\label{sec-demo-translation}

Soit $y[n] = x[n - m_0H]$ avec $m_0$ entier. D'après (2),
\begin{equation}
\begin{split}
  X_y[m,k]
  &= \sum_{n=0}^{N-1} x[n + mH - m_0H]\, w[n]\, e^{-j2\pi kn/\Nfft} \\
  &= \sum_{n=0}^{N-1} x[n + (m-m_0)H]\, w[n]\, e^{-j2\pi kn/\Nfft}
  = X_x[m-m_0, k],
\end{split}
  \label{eq-translation}
\end{equation}
donc $S_y[m,k] = S_x[m-m_0,k]$. Le résultat est exact et vaut même pour la phase, parce que
chaque trame est analysée dans son propre repère temporel. Si $n_0 = m_0H + r$ avec
$0 < r < H$, la trame $m$ de $y$ analyse un segment de $x$ qui commence $r$ échantillons
avant la trame $m-m_0$ : en rapportant chaque trame de $y$ à la trame de $x$ la plus proche,
l'écart est d'au plus $H/2 = N/4$ échantillons, et~\eqref{eq-translation} n'est plus
qu'approchée.

\subsection{Propriété (2) : un filtrage court ajoute un biais \texorpdfstring{$g(k)$}{g(k)}}
\label{sec-demo-filtrage}

Soit $y = \alpha\,(h * x)$ avec $\alpha > 0$ et $h$ causal de longueur $L_h$. On prolonge
$w$ par zéro hors de $[0, N-1]$. En reportant $y[n+mH] = \alpha\sum_l h[l]\,x[n+mH-l]$
dans (2), puis en posant $n' = n - l$ :
\begin{equation}
  X_y[m,k]
  = \alpha \sum_{l=0}^{L_h-1} h[l]\, e^{-j2\pi kl/\Nfft}
     \sum_{n'} x[n'+mH]\, w[n'+l]\, e^{-j2\pi kn'/\Nfft} .
  \label{eq-filtrage-exact}
\end{equation}
Si la fenêtre ne variait pas sur la durée de $h$, on pourrait remplacer $w[n'+l]$ par $w[n']$
et la somme intérieure deviendrait $X_x[m,k]$. En écrivant
$w[n'+l] = w[n'] + (w[n'+l] - w[n'])$, on obtient
$X_y[m,k] = \alpha\,\hhat(f_k)\, X_x[m,k] + \varepsilon[m,k]$ avec
\[
  \varepsilon[m,k] = \alpha \sum_{l} h[l]\, e^{-j2\pi kl/\Nfft}
  \sum_{n'} x[n'+mH]\,\bigl(w[n'+l] - w[n']\bigr)\, e^{-j2\pi kn'/\Nfft} .
\]
Pour la fenêtre de Hann $w[n] = \tfrac12\bigl(1 - \cos(2\pi n/N)\bigr)$, la pente est au
plus $\pi/N$, et la fenêtre est nulle à ses deux bords, donc continue une fois prolongée par
zéro. Par conséquent $|w[n'+l] - w[n']| \le \pi l/N$ pour tout $n'$, bords compris, et
\[
  |\varepsilon[m,k]| \;\le\; \alpha\,\frac{\pi}{N}
  \sum_{l} l\,|h[l]| \sum_{n'=-L_h}^{N-1} |x[n'+mH]| .
\]
Rapportée à $|X_x[m,k]|$, l'erreur est d'ordre $L_h/N$. Avec une fenêtre rectangulaire,
discontinue à ses bords, elle serait d'ordre 1 : la continuité de la fenêtre de Hann est ici
essentielle. Le module au carré puis le logarithme donnent la relation (3).

\subsection{Invariance des maxima au gain et condition (5)}
\label{sec-demo-maxima}

\paragraph{Gain.} Si $y = \alpha x$, alors $S_y = \alpha^2 S_x$, et $s \mapsto \alpha^2 s$ est
strictement croissante : pour tout voisinage $V$,
$S_y[m,k] = \max_{V(m,k)} S_y \iff S_x[m,k] = \max_{V(m,k)} S_x$.

\paragraph{Filtrage.} Soit $(m,k_0)$ un pic de $S_x$, d'avance $\mu$ sur son concurrent le
plus fort dans $V^* = V(m,k_0) \setminus \{(m,k_0)\}$. D'après (3), pour tout
$(m',k') \in V^*$,
\[
  S_{y,\mathrm{dB}}[m,k_0] - S_{y,\mathrm{dB}}[m',k']
  = \underbrace{S_{x,\mathrm{dB}}[m,k_0] - S_{x,\mathrm{dB}}[m',k']}_{\ge\, \mu}
  + \underbrace{g(k_0) - g(k')}_{\ge\, -\Delta g(V)}
  \;\ge\; \mu - \Delta g(V) ,
\]
qui est strictement positif sous la condition (5) : le pic reste le maximum de son voisinage.
La condition est suffisante, pas nécessaire : un pic peut survivre avec $\Delta g(V) \ge \mu$
si son concurrent le plus proche en valeur n'est pas celui que le filtre favorise.

\subsection{Score d'un mauvais morceau et seuil de Gumbel}
\label{sec-annexe-scores}

\paragraph{Modèle uniforme.} Si les clés étaient réparties uniformément sur les $M$ clés
possibles, chacune des $\nq$ clés de la requête trouverait en moyenne $N_c/M$ entrées d'un
morceau $c$ qui en compte $N_c$. Le nombre d'appariements fortuits avec ce morceau suit
alors une loi de Poisson de moyenne $\lambda_c = \nq N_c/M$, répartie sur les
$D \approx (L + \Tq)\Fe/H$ cases de son histogramme des $\delta$ (environ 550 pour des
extraits de \qty{30}{\second}). Chaque case reçoit une loi de Poisson de moyenne
$\lambda = \lambda_c/D$, et
\begin{equation}
  \Prob(\text{fausse détection}) \;\le\; (N_b - 1)\, D\; \Prob\bigl(\mathrm{Poisson}(\lambda) \ge s\bigr) .
  \label{eq-plancher}
\end{equation}
Avec $\nq = 1250$, $N_c \approx \num{6800}$ et $\gf = 1$, $\lambda_c \approx \num{1.1}$, et
$s = 3$ suffit à rester sous \qty{1}{\percent}.

\paragraph{Alphabet effectif.} Si $\pi_u$ est la proportion des entrées de l'index qui
portent la clé $u$, et si les clés de la requête suivent la même loi, la probabilité qu'une
clé de la requête tombe sur une entrée donnée vaut $\sum_u \pi_u^2$ au lieu de $1/M$. Le
modèle reste valable en remplaçant $M$ par $\Meff = 1/\sum_u \pi_u^2 \le M$, avec égalité
seulement si les clés sont uniformes. La contrainte $|\dk| \le 60$ et les bandes
\qtyrange{200}{4000}{\hertz} réduisent déjà l'alphabet accessible à environ \num{1.8e6} clés,
quatre fois moins que $M$.

\paragraph{Notes tenues.} Si la requête contient la clé $u$ à $L$ ancres consécutives
$m, \dots, m+L-1$ et un mauvais morceau à $L'$ ancres consécutives $r, \dots, r+L'-1$, les
$LL'$ appariements donnent $\delta = (r - m) + (i - j)$, et la case $\delta = r - m$ en reçoit
$\min(L, L')$ : des votes cohérents dans le temps, comme pour un bon morceau.

\paragraph{Seuil de Gumbel.} La loi de Gumbel a pour fonction de répartition
$\Prob(S \le s) = \exp\bigl(-e^{-(s-\mu_G)/\beta}\bigr)$. Par la méthode des moments,
\[
  \beta = \frac{\sqrt{6}}{\pi}\,\sigma_s , \qquad \mu_G = \bar{s} - \gamma_E\,\beta ,
  \qquad \gamma_E \approx \num{0.5772},
\]
où $\bar{s}$ et $\sigma_s$ sont la moyenne et l'écart-type des scores de calibration. Pour un
\TFD{} visé $\alpha$, $\Prob(S > \sthr) = \alpha$ donne
$\sthr = \mu_G - \beta \ln\bigl(-\ln(1 - \alpha)\bigr) \approx \mu_G + \num{4.6}\,\beta$ pour
$\alpha = \qty{1}{\percent}$. Passer de $N_b$ à $N_b'$ morceaux décale la loi de
$\beta\ln(N_b'/N_b)$. Mesures (\cref{tab-gumbel}) : le \TFD{} réel est 4 à 5 fois plus élevé
que le \TFD{} visé.

\begin{table}[h]
  \centering\small
  \caption{Seuil de Gumbel et \TFD{} réel mesuré sur les négatives de test des 24 conditions
  (\num{1200} scores).}
  \label{tab-gumbel}
  \begin{tabular}{@{}lcccccc@{}}
    \toprule
    Calibration & $\bar{s}$ & $\sigma_s$ & $\sthr$ ($\alpha = \qty{1}{\percent}$) & \TFD{} réel & $\sthr$ ($\alpha = \qty{0.1}{\percent}$) & \TFD{} réel \\
    \midrule
    50 scores, sans dégradation & \num{7.86} & \num{2.71} & \num{16.4} & \qty{5.0}{\percent} & \num{21.2} & \qty{1.6}{\percent} \\
    \num{1200} scores, toutes conditions & \num{8.09} & \num{3.01} & \num{17.5} & \qty{4.3}{\percent} & \num{22.9} & \qty{1.1}{\percent} \\
    \bottomrule
  \end{tabular}
\end{table}

% ---------------------------------------------------------------------
\section{Notations}
\label{sec-notations}
% ---------------------------------------------------------------------

Les notations sont celles de l'énoncé ; celles que nous ajoutons sont signalées par un
astérisque.

\begin{center}\small
\begin{tabularx}{\linewidth}{@{}lX@{}}
  \toprule
  Symbole & Signification \\
  \midrule
  $x(t)$, $y(t)$ & morceau d'origine, extrait capté \\
  $\alpha$, $h$, $t_0$, $b(t)$ & gain, réponse impulsionnelle haut-parleur $\to$ salle $\to$ micro, décalage, bruit additif (1) \\
  $\Fe$ & fréquence d'échantillonnage \\
  $N$, $H$, $\Nfft$ & longueur de fenêtre, avance entre trames, points de la TFD \\
  $K = \Nfft/2 + 1$ & nombre de canaux \\
  $S[m,k]$, $\SdB[m,k]$ & spectrogramme (trame $m$, canal $k$), en décibels (2) \\
  $t_m = mH/\Fe$, $f_k = k\Fe/\Nfft$ & instant de la trame $m$, fréquence du canal $k$ \\
  $\Tw$, $\Delta f$, $\kw$ & résolution temporelle, résolution fréquentielle, constante de la fenêtre \\
  $g(k)$ & biais introduit par le gain et le filtrage (3) \\
  $V$, $\mv$, $\kv$ & voisinage de $(2\mv+1)$ trames sur $(2\kv+1)$ canaux \\
  $\mu$, $\Delta g(V)$ & avance d'un pic sur ses voisins, variation de $g$ sur le voisinage (5) \\
  $\rho$ & densité de pics (pics/s) \\
  $p$ & probabilité qu'un pic survive à la dégradation \\
  $F$ & nombre de cibles par ancre \\
  $u$, $Q_f$, $\gf$ & clé, classe de fréquence, nombre de canaux par classe (6) \\
  $\dm$, $\dk$ & écarts en trames et en canaux entre ancre et cible \\
  $M = \Nf^2\Nt$ & nombre de clés possibles, $\Nf$ classes de fréquence, $\Nt$ valeurs de $\dm$ \\
  $\Tq$ & durée de la requête \\
  $\nq$, $\Nu$ & nombre de clés de la requête, nombre d'entrées de l'index \\
  $\delta = m_r - m_q$ & décalage entre instants d'ancre dans la référence et dans la requête (7) \\
  $s$, $\sthr$ & score (maximum de l'histogramme des $\delta$), seuil de décision \\
  $N^+$, $N^-$ & nombres de requêtes positives et négatives \\
  \TBD{}, \TFD{} & taux de bonne détection, taux de fausse détection \\
  $\RSB_{\mathrm{dB}}$, $P_x$, $P_b$ & rapport signal à bruit, puissances du signal et du bruit (8) \\
  $L_h$* & longueur de la réponse impulsionnelle $h$ \\
  $N_b$*, $D$* & nombre de morceaux de la base, nombre de cases de l'histogramme des $\delta$ \\
  $\Meff$* & alphabet effectif $1/\sum_u \pi_u^2$ \\
  $\mu_G$*, $\beta$* & paramètres de la loi de Gumbel des scores négatifs \\
  \bottomrule
\end{tabularx}
\end{center}

\section{Liste des morceaux}
\label{sec-liste-morceaux}

Les 200 morceaux utilisés, tous extraits de \texttt{fma\_small}~\cite{fma2017} (extraits de
\qty{30}{\second}, licences Creative Commons). Le fichier complet, avec la licence et la durée
de chaque morceau, est \texttt{liste\_morceaux.csv}. Les titres trop longs sont abrégés, et
quelques caractères non latins ont été translittérés.

% Généré automatiquement depuis liste_morceaux.csv : ne pas modifier à la main.
\begin{footnotesize}
\begin{longtable}{@{}rlll@{}}
\toprule Identifiant & Artiste & Titre & Genre \\ \midrule \endhead
\bottomrule \endfoot
\multicolumn{4}{@{}l}{\textbf{Morceaux de la base (100)}} \\ \midrule
26652 & Graphiqs Groove & Bleuacide & Electronic \\
48439 & Bass Assassin & Panic & Electronic \\
55808 & Blasterhead & Killbots & Electronic \\
67638 & Southman & Wake Up & Electronic \\
69732 & Godblesscomputers & Natürlich & Electronic \\
80793 & Ostin & Spirit & Electronic \\
84198 & Lee Rosevere & Illuminations & Electronic \\
84484 & BLEO & Tart (Pts 1 \& 2) feat. KeFF & Electronic \\
105674 & Gilo & Outro & Electronic \\
107427 & UncleBibby & Brutish & Electronic \\
108415 & UncleBibby & AstroTurf & Electronic \\
142079 & Mike B. Fort & OOOTHERSIDE - Crystal Air & Electronic \\
143046 & Mystery Mammal & Apollo & Electronic \\
22150 & Jeff McLeod & The Horror Approaches & Experimental \\
29752 & Mister Jan & Row Row Row Your Boat (Mister Jan) & Experimental \\
34994 & 6 or 7 & Dear Listener & Experimental \\
39278 & U Can Unlearn Guitar & Chicago, Etc. & Experimental \\
48861 & NaDa BaBa & Skull Flower & Experimental \\
49064 & Justice Yeldham & Track 06 & Experimental \\
52861 & Ilya Monosov & For Jojo Takayanagi & Experimental \\
58054 & Trance And The Arcade & Without Resistance: Nature's Ni\ldots & Experimental \\
68539 & Benboncan & Foyer Crowd & Experimental \\
74392 & Thicker Than Water & Joy to the World & Experimental \\
116514 & SULCUS & R.A.H.M.A.V.N. [221] & Experimental \\
120782 & Teleg. & things are not as they seem & Experimental \\
126521 & L o W \& CROCODILE TEARS & Hymne de apophia & Experimental \\
10443 & Big Blood & Beast (1 Beast) & Folk \\
10993 & Big Blood & Oh Country (Skin \&amp; Bones) & Folk \\
29971 & Amy Sayer & wild tape & Folk \\
47663 & Onda bidón & me aborrezco cada vez que llego\ldots & Folk \\
67016 & Michael Chapman & Kodak Ghosts & Folk \\
88870 & Josh Woodward & Methylchloroisothiazolinone (No\ldots & Folk \\
91443 & Cal Folger Day & Easy Sometimes & Folk \\
108299 & Carter Behnke & Intro & Folk \\
116873 & Itasca & After Dawn & Folk \\
137907 & Matús Novanský & Sally & Folk \\
138416 & Derek Clegg & Egypt, Maybe Rome & Folk \\
139110 & Anthony Pasquarosa & What Makes A Man? & Folk \\
148211 & Josh Armistead & Summertime & Folk \\
11770 & 6th Sense & Future Music Instrumental     \{\ldots & Hip-Hop \\
11775 & 6th Sense & We Don't Stop Instrumental\ldots & Hip-Hop \\
11861 & 6th Sense \&amp; Mick Boogie & Nike Boots Freestyle & Hip-Hop \\
13768 & Majik Most & Just a Feelin' & Hip-Hop \\
14542 & Opio & Due Damage & Hip-Hop \\
19438 & Just Plain Ant & Stumble & Hip-Hop \\
26014 & Comfort Fit & Whenever You Wanna Call Me & Hip-Hop \\
30043 & Cullah & Soul Hip Hop & Hip-Hop \\
72135 & C-Doc & Life Slow Down (Instrumental) & Hip-Hop \\
97037 & CuzOH! Black Feat. Wordsmith & Stronger Living Feat Wordsmith & Hip-Hop \\
108838 & Tha Silent Partner & On Harry Love Drums & Hip-Hop \\
149689 & K. Sparks & Purpose & Hip-Hop \\
150288 & The Impossebulls & Thank U ('92) & Hip-Hop \\
63291 & I Was A Teenage Strangler Sound\ldots & March of the Teenagers pt. 2 & Instrumental \\
86487 & Andrea Carri & Sunshine After Rain & Instrumental \\
92953 & Muhmood & Organ Etude & Instrumental \\
119898 & Blue Dot Sessions & Ferry Landing & Instrumental \\
120298 & Cory Gray & Attic & Instrumental \\
120315 & Cory Gray & Life Size Head & Instrumental \\
125659 & Blue Dot Sessions & Demalion - Theme & Instrumental \\
125681 & Blue Dot Sessions & Eventual Victory & Instrumental \\
125820 & Blue Dot Sessions & Master & Instrumental \\
129042 & Blue Dot Sessions & Third in Line & Instrumental \\
133942 & Raul Diaz Palomar & Plan II & Instrumental \\
141282 & Kevin Hartnell & Rogue Planet & Instrumental \\
6360 & Orkestra Keyif & Sabah Sabah Seyredelim Yaliyi & International \\
6366 & Zikrayat & Taheya & International \\
36272 & Demiran Cerimovic and His Orkestar & Competition Rehearsal, Song 1 & International \\
36273 & Demiran Cerimovic and His Orkestar & Competition Rehearsal, Song 2 & International \\
41018 & Unity Vibration & Au royaume de la money & International \\
54468 & DuBoLoGy & Dub Action & International \\
59718 & Black Sea Hotel & Spavamise & International \\
60865 & Gogofski & Medley: Prishtinka, Ramo-Ramo,\ldots & International \\
62589 & Tchakare Kanyembe & Tchakare Kanyembe track 01 & International \\
73761 & Dad Rocks! & Weapons & International \\
114940 & Zingeray! & Di fishelekh in vaser & International \\
117942 & Volfoniq & L'amour est enfant dub & International \\
22480 & Ergo Phizmiz & Celia & Pop \\
34511 & Symfoniorkestern & Vilse (Slutet Cover) & Pop \\
42911 & Phemale & MEDICAL SEX & Pop \\
70425 & Supercute & Assked & Pop \\
90592 & Pilesar & Bullets Have No Effect & Pop \\
120105 & Lanark & Happy Boccherini & Pop \\
121740 & Frozen Houses & finished enconomic survivor & Pop \\
131452 & Mom Jeans & B06 Say We're Sweethearts Again & Pop \\
131454 & Mom Jeans & B08 Mario & Pop \\
135339 & Caleb Lemond & I Think From Now On & Pop \\
148286 & LoveyDove & Beppe of Italy & Pop \\
148288 & LoveyDove & Deep Down Inc. & Pop \\
42851 & Waylon Thornton and the Heavy H\ldots & Manson Halloween & Rock \\
44918 & Apache Dropout & Teenager & Rock \\
55718 & et\_ & Kopeika & Rock \\
68821 & Zulus & Death in the Current & Rock \\
69781 & Home Blitz & Secret Wave & Rock \\
96400 & Insane Ride & Save Your Tears & Rock \\
110634 & Eidetic Seeing & No Pilot & Rock \\
119257 & The Wades & like love & Rock \\
122397 & Gorilla Chief & Cocaine Girl & Rock \\
125001 & Let Me Crazy & 899 Seconds & Rock \\
126300 & Andy G. Cohen & Scramby Eggs & Rock \\
142511 & Miët & Ratatatatouille & Rock \\
\midrule
\multicolumn{4}{@{}l}{\textbf{Morceaux hors base, requêtes négatives (100)}} \\ \midrule
23016 & The Lights Galaxia & While She Sleeps (Morning Edit) & Electronic \\
48465 & Sycamore Drive & Ocean Breeze & Electronic \\
52409 & Buskerdroid & I Want You & Electronic \\
52630 & Jason Shaw & Quirky Jerk & Electronic \\
65063 & Sycamore Drive & Kicks & Electronic \\
66390 & krackatoa & Komodo Dodo & Electronic \\
69211 & JMD v Kube & Interference (Psychadelik Pedes\ldots & Electronic \\
72067 & Decktonic & Africa & Electronic \\
85831 & Frame & Flames (The Beginning) & Electronic \\
87070 & (F.E.V.E.R.) & Argonaut Into Deep Space (SFI r\ldots & Electronic \\
92366 & Out Of Orion (Ox3) & Horror House & Electronic \\
126899 & SPCZ & Kuma & Electronic \\
146458 & Langax & The Cavern & Electronic \\
38891 & File Under Toner & justified & Experimental \\
44781 & Logan Whitehurst and The Jr. Sc\ldots & (She's) Ten Feet Tall & Experimental \\
49068 & Justice Yeldham & Knobranaut & Experimental \\
51158 & Base Cleft & Carousel of Bent Facts & Experimental \\
56015 & Tone Poets & Surface & Experimental \\
63161 & Lucky Dragons & 13 & Experimental \\
73821 & Miquel Parera Jaques & nxVacuity002-2012-12-12\_04\_00\_16 & Experimental \\
110611 & Jeff Kolar & Happy Birthday To You & Experimental \\
112668 & ODOR BABY & 45 & Experimental \\
112788 & ODOR BABY & BROKEN NEWS & Experimental \\
124870 & Ross Oldenburg & Five-O (Improvisation for Five\ldots & Experimental \\
133294 & AG Davis & Jazz Got Fuck'd & Experimental \\
145043 & Hermit & Rock'n'Roll Hara-Kiri & Experimental \\
12058 & Big Blood & Land Upon Mud I & Folk \\
12060 & Big Blood & Heart of Glass & Folk \\
12064 & Big Blood & Beatle Bones 'N' Smokin' Stones & Folk \\
13930 & Foghorn Stringband & John Henry & Folk \\
41192 & Weyes Blood & Storms That Breed & Folk \\
63253 & Big Blood & Sirens Knell & Folk \\
69765 & Plusplus & Echoes & Folk \\
89477 & Wellington Sea Shanty Society & Great Open Sea & Folk \\
130955 & Albin Andersson & Even If It's True & Folk \\
135044 & Anga & Le diable & Folk \\
140791 & NaDa BaBa & Sam's Song & Folk \\
144475 & Nick Jaina & The Dinner Party (Instrumental) & Folk \\
145551 & This Mess is Mine & Thief & Folk \\
22091 & ETHX & Wondering Why (Featuring Nate M\ldots & Hip-Hop \\
43867 & Chicken Jones & one last try & Hip-Hop \\
56498 & Opio & World of Pain & Hip-Hop \\
63755 & Cullah & A Dangerous Game & Hip-Hop \\
63803 & Lushlife & Gymnopedie (feat. Shad) & Hip-Hop \\
71690 & C-Doc & Drumm Machine (12 Inch Mixx) (F\ldots & Hip-Hop \\
75372 & Kellee Maize & 2012 & Hip-Hop \\
75425 & Kellee Maize & Eleven 11 Remix & Hip-Hop \\
97697 & Tickle & Stand Up & Hip-Hop \\
116390 & Asthmatic Astronaut & Jazz soul & Hip-Hop \\
129914 & Dollar Boyz & King Of Da City & Hip-Hop \\
132962 & Cockmasta Flex & Maricones Calientes (remastered) & Hip-Hop \\
140623 & The Impossebulls & The More Things Change & Hip-Hop \\
54153 & Mark Neil & 05 TheLastDodo TheHappyAndroid & Instrumental \\
75435 & Cherly KaCherly & Time is running out. I'm sure. & Instrumental \\
114212 & Squire Tuck & Old Fashioned Lullaby & Instrumental \\
114534 & Blue Dot Sessions & Diner Gothic - Dasher & Instrumental \\
116177 & Squire Tuck & A Mexican Love Affair & Instrumental \\
121651 & happiness in aeroplanes & Her History of Water & Instrumental \\
123761 & Hawkin & Week 1 - rain & Instrumental \\
123764 & Hawkin & Week 4 - clarity & Instrumental \\
126411 & Blue Dot Sessions & Chromium Blush & Instrumental \\
142641 & TRG Banks & Goodbye Machineryhead & Instrumental \\
144454 & TRG Banks & Father Green & Instrumental \\
148132 & Ryan Cullinane & Upside of Looking Down Instrume\ldots & Instrumental \\
1688 & The Sounds of Taraab & Sishituki & International \\
44822 & Black Bear Combo & Watchyerstep & International \\
54469 & DuBoLoGy & Walk In Forest & International \\
54482 & Friglob & Love music & International \\
60753 & Niva & Snosti ti Dojdov & International \\
80753 & Gogofski & Medley: Sadila Jana, Davaj Me & International \\
99362 & Svitanya & Vet Me The & International \\
108867 & Black Sea Hotel & Malka Moma Dvora Mela & International \\
129926 & Radio Jarocho and Zenen Zeferin\ldots & Vieja & International \\
131941 & Choban Elektrik & Unknown Cocek Tune & International \\
150080 & DubRaJah & Deep Karnataka & International \\
152253 & Panda Dub & Milky Way & International \\
28485 & The Pop Winds & Good Person & Pop \\
32334 & Plushgoolash & Brown Rainbows & Pop \\
37369 & Útidúr & Grasping for Thoughts & Pop \\
59445 & Phemale & This Place & Pop \\
61670 & Gosprom & Japan Punk & Pop \\
61736 & Soft Serve & Margaritaville & Pop \\
68353 & Caroline Gorman & Hiding in the Dark & Pop \\
91160 & Monroeville Music Center & Invaders From Within & Pop \\
127184 & Wintermitts & I Saw A Ghost & Pop \\
129391 & A la Mode & Mountaintop & Pop \\
134385 & Scott Holmes & Moving Forward & Pop \\
138038 & Massage & Breaking Up & Pop \\
13556 & STEEP & Let Me Go & Rock \\
69723 & ROGER ROBERT & Dracula & Rock \\
93986 & Bugs and Rats & Irish Clouds & Rock \\
106277 & Night Rally & A Birthday Party & Rock \\
106570 & Apache Tomcat & Texas Ostrich Farm & Rock \\
108502 & Voodoo Puppets & Slender Man & Rock \\
110985 & Antethic & Cheliuskin & Rock \\
129973 & Pipe Choir & Savior & Rock \\
136707 & Tojamura 7 & 100ft Wave & Rock \\
138060 & Every Now and Every Then & One Knife & Rock \\
146478 & Norwegian Blue Records & SPARKLE BLOOD - Denim Convention & Rock \\
148537 & E A Terra Nunca Me Pareceu Tão\ldots & Medo de Morrer & Rock \\
\midrule
\end{longtable}
\end{footnotesize}

